% A5 Fraction guide. Build with: xelatex Guide.tex
\documentclass[11pt,a4paper]{article}

\usepackage[margin=2.2cm]{geometry}
\usepackage{fontspec}
\usepackage{xeCJK}
\setmainfont{TeX Gyre Heros}
\setmonofont{DejaVu Sans Mono}[Scale=0.88]
\setCJKmainfont{WenQuanYi Zen Hei}[AutoFakeBold=2.5]
\setCJKmonofont{WenQuanYi Zen Hei}
\xeCJKsetup{CJKecglue={\hskip 0.2em plus 0.1em}}

\usepackage{xcolor}
\usepackage{titlesec}
\usepackage{booktabs}
\usepackage{tabularx}
\usepackage{array}
\usepackage{enumitem}
\usepackage{listings}
\usepackage[most]{tcolorbox}
\usepackage[hidelinks]{hyperref}

\definecolor{accent}{HTML}{1F5FAD}
\definecolor{codebg}{HTML}{F5F7FA}
\definecolor{donebg}{HTML}{EEF6EE}
\definecolor{donefr}{HTML}{3C8D40}
\definecolor{warnbg}{HTML}{FFF6E5}
\definecolor{warnfr}{HTML}{D08A00}

\titleformat{\section}{\Large\bfseries\color{accent}}{}{0pt}{}
\titleformat{\subsection}{\large\bfseries}{}{0pt}{}
\titlespacing*{\section}{0pt}{2.2ex plus .5ex}{1ex}
\titlespacing*{\subsection}{0pt}{1.8ex plus .4ex}{0.8ex}
\setlength{\parindent}{0pt}
\setlength{\parskip}{0.7ex}
\linespread{1.15}
\setlist{nosep, leftmargin=1.6em, topsep=0.3ex}
\renewcommand{\arraystretch}{1.25}

\lstset{
  language=C++, basicstyle=\ttfamily\small, columns=fullflexible, keepspaces=true,
  backgroundcolor=\color{codebg}, frame=single, rulecolor=\color{codebg},
  framesep=4pt, xleftmargin=4pt, xrightmargin=4pt, aboveskip=0.8ex, belowskip=0.8ex,
  keywordstyle=\color{accent}\bfseries, commentstyle=\color{gray}, stringstyle=\color{donefr},
  showstringspaces=false, morekeywords={explicit}
}

\newcommand{\code}[1]{\texttt{#1}}
\newcommand{\idea}{\textbf{思路：}}
\newcommand{\todo}{\textbf{你来写：}}

\newtcolorbox{done}{enhanced, before skip=2.2ex, colback=donebg, colframe=donefr, boxrule=0.6pt, arc=2pt,
  left=6pt, right=6pt, top=8pt, bottom=4pt, fontupper=\small,
  title={完成标志}, coltitle=white, fonttitle=\small\bfseries,
  attach boxed title to top left={xshift=6pt, yshift=-2mm},
  boxed title style={colback=donefr, colframe=donefr, arc=1pt, boxrule=0pt}}
\newtcolorbox{tip}[1][小知识]{enhanced, before skip=2.2ex, colback=warnbg, colframe=warnfr, boxrule=0.6pt, arc=2pt,
  left=6pt, right=6pt, top=8pt, bottom=4pt, fontupper=\small,
  title={#1}, coltitle=white, fonttitle=\small\bfseries,
  attach boxed title to top left={xshift=6pt, yshift=-2mm},
  boxed title style={colback=warnfr, colframe=warnfr, arc=1pt, boxrule=0pt}}

\begin{document}
\sloppy
\setlength{\emergencystretch}{3em}

\begin{center}
  {\LARGE\bfseries A5 Fraction：动手指南}\\[0.8ex]
  {\normalsize Assignment 005 · 分数类}
\end{center}

\textbf{目标：}写一个表示分数的类 \code{Fraction}，让它像内置数字一样能加减乘除、比较大小、打印和读入。

本文是一条推荐路线，不是唯一写法。每一步写完都要编译、跑一下“完成标志”里的例子，通过了再做下一步。

\section{一、题目要求对照}

题目列了 9 项，先看清每一项在 C++ 里对应什么写法：

\begin{center}
\small
\begin{tabularx}{\linewidth}{@{}>{\raggedright\arraybackslash}p{3.3cm}>{\raggedright\arraybackslash}X>{\raggedright\arraybackslash}p{4.6cm}@{}}
\toprule
题目要求 & C++ 写法 & 例子 \\
\midrule
Default ctor & \code{Fraction()} & \code{Fraction z;} 得到 \code{0/1} \\
两个整数的构造 & \code{Fraction(long long n, long long d)} & \code{Fraction(2, -4)} 得到 \code{-1/2} \\
Copy ctor & \code{Fraction(const Fraction\& other)} & \code{Fraction b(a);} \\
\code{+ - * /} & \code{operator+} 等，共 4 个 & \code{a + b} \\
\code{< <= == != >= >} & \code{operator<} 等，共 6 个 & \code{a < b} \\
转换为 double & \code{explicit operator double() const} & \code{static\_cast<double>(a)} \\
To string & \code{std::string toString() const} & \code{"2/3"} \\
流的插入和提取 & \code{operator<<}、\code{operator>>} & \code{cout << a;}\ \code{cin >> a;} \\
小数字符串转换 & \code{static Fraction fromDecimal(const std::string\&)} & \code{"1.414"} 得到 \code{707/500} \\
\bottomrule
\end{tabularx}
\end{center}

“运算符重载”（operator overloading）就是：给自己的类写一个名字叫 \code{operator+} 的函数，之后写 \code{a + b} 时，编译器会自动调用 \code{operator+(a, b)}。它和普通函数没有本质区别，只是名字和调用方式特殊。

\section{二、先定三条约定}

\subsection{1. 永远保存最简形式}

规定：分母永远大于 0，分子分母互质，0 写成 \code{0/1}。例如 \code{2/-4} 存成 \code{-1/2}。

这样做的好处：
\begin{itemize}
  \item 同一个数只有一种写法，\code{==} 只要比较分子和分母是否都相等；
  \item 分母为正，比较大小时交叉相乘不会变号（第 4 步会用到）；
  \item 输出统一，测试时容易写期望结果。
\end{itemize}

关键技巧：\textbf{所有分数都通过“两个整数的构造函数”生成，化简只写在这一个地方。}后面的加减乘除只要算出新的分子分母再交给构造函数，结果自动是最简的。

\subsection{2. 用 long long 存分子分母}

加法要算 \code{a*d + b*c}，乘法会让数字变大；\code{"1.414"} 也要先变成 \code{1414/1000}。用 \code{long long}（约 $\pm 9.2\times10^{18}$）比 \code{int} 安全得多。超大数仍可能溢出，本作业不处理，在 README 里写一句即可。

\subsection{3. 出错就抛异常}

三种错误：分母为 0、除以 0、字符串不是合法小数。构造函数不能像 A4 的 \code{Date::parse} 那样返回 \code{bool}，所以统一抛出标准异常 \code{std::invalid\_argument}（头文件 \code{<stdexcept>}）。只需要会两句：

\begin{lstlisting}
throw std::invalid_argument("denominator is zero");   // where the error happens

try {                                                  // where you test it
    Fraction bad(1, 0);
} catch (const std::invalid_argument& e) {
    std::cout << "rejected: " << e.what() << '\n';
}
\end{lstlisting}

只有一个例外：\code{operator>>} 读到坏数据时不抛异常，而是让流进入失败状态（第 6 步讲）。

\section{三、代码怎么分}

\begin{center}
\small
\begin{tabular}{@{}ll@{}}
\toprule
文件 & 内容 \\
\midrule
\code{Fraction.h} & 类的声明和所有运算符的声明 \\
\code{Fraction.cpp} & 全部实现 \\
\code{main.cpp} & 测试：逐项调用并打印 PASS / FAIL \\
\bottomrule
\end{tabular}
\end{center}

下面是推荐的头文件骨架。它只是“菜单”，函数体要你在 \code{Fraction.cpp} 里自己写：

\begin{lstlisting}
#ifndef FRACTION_H
#define FRACTION_H

#include <iosfwd>
#include <string>

class Fraction {
public:
    Fraction();                                       // 0/1
    Fraction(long long numerator, long long denominator);
    Fraction(const Fraction& other);
    Fraction& operator=(const Fraction& other) = default;

    static Fraction fromDecimal(const std::string& text);   // "1.414"
    std::string toString() const;                            // "2/3"
    explicit operator double() const;

    long long numerator() const { return num; }
    long long denominator() const { return den; }

private:
    long long num;   // keeps the sign
    long long den;   // always > 0
};

Fraction operator+(const Fraction& a, const Fraction& b);
// ... operator- * /  and  < <= == != >= >  (same parameter list)
std::ostream& operator<<(std::ostream& os, const Fraction& f);
std::istream& operator>>(std::istream& is, Fraction& f);

#endif
\end{lstlisting}

\begin{tip}[为什么运算符写在类外面]
\code{a + b} 两边地位相同，写成类外的普通函数（非成员函数）最自然。它们通过公有的 \code{numerator()}、\code{denominator()} 读取数据，不需要 \code{friend}。成员函数后面的 \code{const} 表示“这个函数不修改对象”。
\end{tip}

\section{四、按六步实现}

\subsection{第 1 步：构造函数与化简}

\idea 化简只写在两个整数的构造函数里，其他构造都很简单。

\todo
\begin{itemize}
  \item 默认构造：分子 0，分母 1。
  \item 两个整数的构造，按顺序做三件事：
  \begin{enumerate}
    \item 分母为 0：抛异常；
    \item 分母为负：分子分母同时取反；
    \item 两者同时除以最大公约数：\code{std::gcd}（头文件 \code{<numeric>}，C++17）。
  \end{enumerate}
  \item 拷贝构造：逐个复制两个成员，用初始化列表 \code{: num(other.num), den(other.den)}。
\end{itemize}

\begin{tip}
\begin{itemize}[leftmargin=1.2em]
  \item \code{std::gcd(0, 5)} 等于 5，所以 \code{0/5} 化简后自然是 \code{0/1}，不用特殊处理。
  \item 拷贝构造在这些时候被调用：\code{Fraction b(a);}、\code{Fraction b = a;}、按值传参、按值返回。
  \item 骨架里的 \code{operator= = default} 不能省：自己写了拷贝构造后，编译器默认生成的赋值运算符被标为“过时”，\code{-Wextra} 会在每个 \code{d = c;} 处报 \code{-Wdeprecated-copy} 警告（已实测）。
\end{itemize}
\end{tip}

\begin{done}
暂时用 \code{numerator()} 和 \code{denominator()} 打印检查：
\code{Fraction(2, -4)} 是 $-1$ 和 $2$；\code{Fraction(0, -5)} 是 $0$ 和 $1$；\code{Fraction()} 是 $0$ 和 $1$；
\code{Fraction(1, 0)} 抛出异常。
\end{done}

\subsection{第 2 步：toString 和 \texttt{<<}}

\idea 先让分数能打印出来，后面每一步都靠它检查结果。

\todo
\begin{itemize}
  \item \code{toString()}：返回 \code{"分子/分母"}，用 \code{std::to\_string} 拼接。整数也写成 \code{3/1}，这样输出和输入格式完全一致。
  \item \code{operator<<}：把 \code{toString()} 的结果写进 \code{os}，然后 \code{return os;}。
\end{itemize}

\begin{tip}
\code{operator<<} 返回 \code{std::ostream\&}，是为了支持链式输出：\code{cout << a << " " << b} 先算 \code{cout << a}，得到的还是 \code{cout}，再接着输出后面的内容。
\end{tip}

\begin{done}
\code{std::cout << Fraction(2, -4) << " " << Fraction() << '\textbackslash n';} 输出 \code{-1/2 0/1}。
\end{done}

\subsection{第 3 步：算术运算符}

\idea 按公式算出新的分子和分母，交给构造函数，化简和除零检查都自动完成。

\begin{center}
\small
\begin{tabular}{@{}ll@{}}
\toprule
运算 & 公式 \\
\midrule
加 & $\dfrac{a}{b} + \dfrac{c}{d} = \dfrac{ad + cb}{bd}$ \\[1.2ex]
减 & $\dfrac{a}{b} - \dfrac{c}{d} = \dfrac{ad - cb}{bd}$ \\[1.2ex]
乘 & $\dfrac{a}{b} \times \dfrac{c}{d} = \dfrac{ac}{bd}$ \\[1.2ex]
除 & $\dfrac{a}{b} \div \dfrac{c}{d} = \dfrac{ad}{bc}$ \\
\bottomrule
\end{tabular}
\end{center}

示范加法，另外三个照着写：

\begin{lstlisting}
Fraction operator+(const Fraction& a, const Fraction& b) {
    return Fraction(a.numerator() * b.denominator() + b.numerator() * a.denominator(),
                    a.denominator() * b.denominator());
}
\end{lstlisting}

除以 0 时，除法公式的新分母 $bc$ 等于 0，构造函数会自己抛异常，\code{operator/} 里不用再写判断。

\begin{done}
\code{1/2 + 1/3 = 5/6}，\code{1/2 - 1/3 = 1/6}，\code{1/2 * 1/3 = 1/6}，\code{(1/2) / (1/3) = 3/2}；
\code{Fraction(1, 2) / Fraction()} 抛出异常。
\end{done}

\subsection{第 4 步：关系运算符}

\idea 只认真实现 \code{<} 和 \code{==}，另外四个用它们组合，不重复写比较逻辑。

\todo
\begin{itemize}
  \item \code{==}：两个数都是最简形式，所以分子相等且分母相等即可。
  \item \code{<}：$\dfrac{a}{b} < \dfrac{c}{d} \iff ad < cb$。因为 $b, d > 0$，两边同乘 $bd$ 不改变不等号方向，这就是约定“分母为正”的用处。
  \item 其余四个：
\end{itemize}

\begin{center}
\small
\begin{tabular}{@{}llll@{}}
\toprule
\code{a != b} & \code{a > b} & \code{a <= b} & \code{a >= b} \\
\midrule
\code{!(a == b)} & \code{b < a} & \code{!(b < a)} & \code{!(a < b)} \\
\bottomrule
\end{tabular}
\end{center}

\begin{done}
对 \code{1/2} 和 \code{1/3}：\code{<}、\code{<=}、\code{==} 为假，\code{!=}、\code{>=}、\code{>} 为真；
\code{Fraction(1, 2) == Fraction(2, 4)} 为真；\code{Fraction(-1, 2) < Fraction(1, 3)} 为真。
\end{done}

\subsection{第 5 步：转换为 double}

\todo 写 \code{explicit operator double() const}。计算时先把分子转成 \code{double} 再除：\code{static\_cast<double>(num) / den}，否则 \code{2 / 3} 是整数除法，结果为 0。

\begin{tip}[为什么一定要加 explicit]
不加 \code{explicit} 时，编译器会在你没注意的地方自动把分数变成小数。实测：\code{f} 为 \code{1/2} 时，\code{f + 1} 能编译通过，得到 \code{double} 类型的 \code{1.5}；\code{f == 0.5} 也能编译通过。分数悄悄丢失了精度，编译器却一声不吭。加上 \code{explicit} 后，必须写 \code{static\_cast<double>(f)} 才能转换。
\end{tip}

\begin{done}
\code{static\_cast<double>(Fraction(2, 3))} 输出 \code{0.666667}；\code{Fraction(1, 4)} 转为 \code{0.25}。
\end{done}

\subsection{第 6 步：输入和小数字符串}

\textbf{(1) \texttt{operator>>}}

\idea 读入格式和 \code{toString()} 一致：\code{分子/分母}。

\todo
\begin{itemize}
  \item 用 \code{is >> n >> slash >> d} 依次读入分子、斜杠和分母（\code{slash} 是 \code{char}）。
  \item 读取成功后检查：\code{slash == '/'} 且 \code{d != 0}，满足就 \code{f = Fraction(n, d);}，否则调用 \code{is.setstate(std::ios::failbit);}。
  \item 最后 \code{return is;}。
\end{itemize}

这里不抛异常，而是让流“失败”，因为大家习惯写 \code{while (cin >> f)}：流失败时循环自然停止。通过公有构造函数给 \code{f} 赋值，也不需要 \code{friend}。

\textbf{(2) \texttt{fromDecimal}}

\idea 把 \code{1.414} 看成 \code{1414 / 1000}，再交给构造函数化简成 \code{707/500}。它是静态成员函数，理由和 A4 的 \code{Date::parse} 一样：调用之前还没有分数对象。

\todo 逐个字符扫描：
\begin{enumerate}
  \item 开头可以有一个 \code{+} 或 \code{-}；
  \item 遇到数字：\code{n = n * 10 + (c - '0')}；如果已经过了小数点，分母 \code{d *= 10}；
  \item 遇到小数点：只允许出现一次；
  \item 其他字符、一个数字都没有、数字超过 18 位（防止 \code{long long} 溢出）：抛异常；
  \item 最后返回 \code{Fraction(带符号的 n, d)}。
\end{enumerate}

\begin{done}
\begin{tabular}{@{}ll@{\hspace{2.5em}}ll@{}}
\code{"1.414"} & \code{707/500} & \code{""} & 抛异常 \\
\code{"-0.5"} & \code{-1/2} & \code{"-"} & 抛异常 \\
\code{"3"} & \code{3/1} & \code{"1.2.3"} & 抛异常 \\
\code{".25"} & \code{1/4} & \code{"abc"}、\code{"1e5"} & 抛异常 \\
\end{tabular}

\smallskip
用 \code{std::istringstream in("3/4 -5/10 7/0");}（头文件 \code{<sstream>}）循环读入：依次读到 \code{3/4}、\code{-1/2}，在 \code{7/0} 处停止，\code{in.fail()} 为真。
\end{done}

\section{五、测试}

在 \code{main.cpp} 里写一个小函数，把“完成标志”里的例子逐条变成检查：

\begin{lstlisting}
void check(bool ok, const std::string& name) {
    std::cout << (ok ? "PASS  " : "FAIL  ") << name << '\n';
}

check(Fraction(1, 2) + Fraction(1, 3) == Fraction(5, 6), "1/2 + 1/3 == 5/6");
check(Fraction::fromDecimal("1.414").toString() == "707/500", "1.414 -> 707/500");
\end{lstlisting}

至少覆盖下面这些情况：

\begin{center}
\small
\begin{tabularx}{\linewidth}{@{}lX@{}}
\toprule
类别 & 用例 \\
\midrule
构造与化简 & 默认构造、负分母、需要约分、分子为 0、分母为 0（异常）、拷贝构造、赋值 \\
算术 & 四种运算各一例，含负数；除以 0（异常） \\
比较 & 六个运算符；相等但写法不同（\code{1/2} 与 \code{2/4}）；负数比较 \\
转换 & 转 double；\code{toString} \\
流 & \code{<<} 输出；\code{>>} 正常读入、遇到 \code{7/0} 和非法字符时失败 \\
小数字符串 & 正数、负数、整数、\code{.25}；空串、只有符号、两个小数点、非法字符（异常） \\
\bottomrule
\end{tabularx}
\end{center}

异常用例用 \code{try / catch} 包起来：没抛异常算 FAIL，抛了算 PASS。编译运行：

\begin{lstlisting}[language=bash]
g++ -std=c++17 -Wall -Wextra Fraction.cpp main.cpp -o fraction
./fraction
\end{lstlisting}

全部 PASS，并且编译没有警告，才算完成。

\section{六、最后提交什么}

\begin{itemize}
  \item \code{code/}：\code{Fraction.h}、\code{Fraction.cpp}、\code{main.cpp}；不放可执行文件和 \code{.o}。
  \item \code{README.md}：环境（系统、编译器版本）、编译和运行命令、你的约定：输出格式 \code{n/d}、输入格式、小数格式、哪些情况抛异常。
  \item \code{report/}：按课程命名的 PDF，包含环境、运行方式和真实的测试结果。
\end{itemize}

把它们放进 \code{Assignment05\_<学号>\_<姓名拼音>/}，打包成 \code{.zip}。

\bigskip
\begin{center}
\textbf{现在只做第 1 步：写出三个构造函数，验证化简和除零检查。通过后再开始第 2 步。}
\end{center}

\end{document}
